%%%%%%%%%%%%%%%%%%%%%
%%%%%%%%%%%%%%%%%%%%%
\section{Mapping proteins into property vectors}
\label{sec:vqfeatures}
%%%%%%%%%%%%%%%%%%%%%
%%%%%%%%%%%%%%%%%%%%%

%For a sequence $\bf{s}$ we can define a vector \textbf{$\mathbf{p}$} as any mapping of $\mathcal{A}^{\ell}$ to $\mathbb{R}^{\ell}$. For example, $\mathbf{p}$ may be provided as a vector of hydrophobicity indices corresponding to amino acids in $\mathbf{s}$. Alternatively, it can be represented as predicted helical propensities as outputted by some predictor of secondary structure. For those sequences with available structures, $\mathbf{p}$ may correspond to a sequence of dihedral angles calculated from the protein structure model of $\mathbf{s}$. 

%%%%%%%%%%%%%%%%%%%%%
%%%%%%%%%%%%%%%%%%%%%
\subsection{Structure based properties}
%%%%%%%%%%%%%%%%%%%%%
%%%%%%%%%%%%%%%%%%%%%

Structure based properties were generated by converting the 3D coordinates of amino acids into backbone angles. These angles can be used to position each amino acid in relation to surrounding atoms in the protein chain.  The usefulness of representing a structure in this manner is that backbone angles are the same regardless of how the overall protein is oriented in space, whereas 3D coordinates can change for the exact same structure simply by rotating or shifting it.  

%%%%%%%%%%%%%%%%%%%%%%%
%%%%%%%%%%%%%%%%%%%%%%%
\begin{figure}[h!]
\begin{centering}

	\begin{subfigure}[b]{0.40\textwidth}	
		\centering
		\includegraphics[width=\textwidth]{Figures/VQ/phi_psi}
		\label{fig:phipsibackbone}
		\caption{$\phi$ and $\psi$ tortion angles}
	\end{subfigure}
	\quad
	\begin{subfigure}[b]{0.40\textwidth}	
		\centering
	     	\includegraphics[height=45mm]{Figures/VQ/kappa_alpha}
		\label{fig:kappaalphabackbone}
		\caption{$\kappa$ and $\alpha$ angles}
	\end{subfigure}

	
\caption[Description of angles derived from protein structures]{The different types of angles that can be measured from protein structures are shown.  \Cref{fig:phipsibackbone} illustrates the backbone torsion angles, $\phi$ and $\psi$, for a $C_{\alpha}$ atom that can be measured between its adjacent nitrogen and $C^{\prime}$ atoms respectively. \Cref{fig:kappaalphabackbone} illustrates the angles that can be measures between successive $C_{\alpha}$ atoms: $C^{i-1}_{\alpha}$ and $C^{i+2}_{\alpha}$ in the case of $\kappa$ angles; and $C^{i-2}_{\alpha}$ and $C^{i+2}_{\alpha}$ in the case of $\alpha$ angles. }
\end{centering}
\label{fig:vqbackbone}
\end{figure}
%%%%%%%%%%%%%%%%%%%%%%%
%%%%%%%%%%%%%%%%%%%%%%%

We utilized four types of backbone angles: $\kappa$, $\alpha$, $\phi$, and $\psi$.  \Cref{fig:vqbackbone} illustrates how these angles are defined using a protein chain. As illustrated by \Cref{fig:phipsibackbone}, $\phi$ and $\psi$ angles represent the torsion angle, or amount of rotation, of the bond between $C_{\alpha}$ side-chain  attached carbon atom and the adjacent nitrogen ($\phi$ angle) and $C^{\prime}$ carbon ($\psi$ angle) atoms. While $\phi$ and $\psi$ represent torsion angles between adjacent atoms in a protein's backbone, $\kappa$ and $\alpha$ angles represent the relationship between more distant $C_{\alpha}$ atoms.  An $\ell$ amino acid length sequence, $\mathbf{s}=s_{1},s_{2}\ldots s_{\ell}$, will have $\ell$ associate $C_{\alpha}$ carbon atoms:   $C_{\alpha}^{1}, C_{\alpha}^{2} \ldots C_{\alpha}^{\ell}$. Using \Cref{fig:kappaalphabackbone} as a reference, if one designates a $C_{\alpha}^{i}$ atom as a vertex (in the angular sense), the $\kappa$ angle represents the angle obtained by extending a line from the $C_{\alpha}^{i}$ atom's position to the $C_{\alpha}^{i-2}$ and $C_{\alpha}^{i+2}$ atoms.  The $\alpha$ angle represents the dihedral angle measured between the $C_{\alpha}^{i-1}$ and $C_{\alpha}^{i+2}$ atoms when two planes are generated using the $C_{\alpha}^{i}$ atom to the $C_{\alpha}^{i+1}$ to define the line where the two planes meet.

Angles $\kappa$, $\alpha$, $\phi$, and $\psi$ were obtained using DSSP \citep{Kabsch1983}.  In addition to generating backbone angles, DSSP also outputs solvent accessibility values which we used as the fifth structure based property.

%%%%%%%%%%%%%%%%%%%%%
%%%%%%%%%%%%%%%%%%%%%
\subsection{Sequence based properties}
%%%%%%%%%%%%%%%%%%%%%
%%%%%%%%%%%%%%%%%%%%%

We generated seven sequence based properties for both the task of categorizing structures into SCOP folds and predicting GO function.  These features were generated in order to represent a biologically relevant property associated with a region of a protein sequence: ($i$) hydrophobicity, calculated using the the Kyte-Doolittle scale \citep{Kyte1982} in a sliding window with $w=11$; ($ii$) flexibility, calculated as predicted B-factors using our previous model \citep{radivojac2004protein}; secondary structure predictions of ($iii$) helix, ($iv$) loop and ($v$) sheet propensities using our in-house predictor; and intrinsic disorder, ($vi$) using the previously published VSL2B model \citep{Peng2006} as well as ($vii$) predictions from the same in-house predictor used for secondary structures. 

\subsection{Tested values of window size and number of centroids}
We tested a range of combinations of values for $n$ (window size), and $m$ (number of centroids) for each property.  For values of $n$, we tested $n\in \{2^{i}:i=2\dots 5\}$.  Similarly, for the number of centroids, $m$, we tested $m\in\{2^{i}:i=6,8,10,12\}$. For all values of $m$ and $n$ we performed k-means clustering using $10^6$ randomly sampled vectors from all data points.

